\section{Adaptive finite-rank constraints on small patches}
\label{sec:patches}

The next construction gives an accurate constraint near any specified
square of logarithmic radius.  Its conclusion allows the radius of the
finite-rank part to vary between terms.  This freedom is essential:
the proposition does not assert an accurate finite-rank truncation at
one prescribed boundary.  We use a regularized version of the nested
filter method in \cite{RepoAreaLaw}.  The optimizing filters need not
commute, so their conjugations must be handled in the order supplied by
the nested squares.

For a center \(z\in\mathbb R^2\) and \(r\geq 0\), write
\[
 X(r)=\{v\in\Lambda:\abs{v-z}_\infty\leq r\},
 \qquad
 \mathcal H_Y=\bigotimes_{v\in Y}\mathbb C^q.
\]
In particular, these are closed square samples, possibly clipped by the
physical boundary.  The Hilbert space of an empty set is \(\mathbb C\).
For a subspace $S\subseteq\mathcal H_X$, the cylinder
$S\otimes\mathcal H_{\Lambda\setminus X}$ has \emph{inside rank} $\dim S$.

\begin{proposition}[Adaptive patch constraint]
\label{prop:patch}
Fix \(B_0>0\) and \(p>0\).  There are constants \(C,c>0\), depending
only on \(q,J,\Delta,B_0,p\), with the following property.  For every
\(L\geq2\), center \(z\in\mathbb R^2\), and
\(0<u\leq B_0\log L\), there is an orthogonal projection \(P\),
supported on \(X(2u)\), such that
\[
 \norm{(\id-P)\Omega}\leq L^{-p}.
\]
There are radii \(u\leq r_1\leq\cdots\leq r_m\leq2u\), sets
\(X_j=X(r_j)\), and unit vectors
\(v_{j\ell}\in\mathcal H_{X_j}\), \(1\leq\ell\leq d_j\), for which
\begin{equation}
 P=\sum_{j=1}^m\sum_{\ell=1}^{d_j}
       \lvert v_{j\ell}\rangle\langle v_{j\ell}\rvert_{X_j}
          \otimes\id_{\Lambda\setminus X_j},
 \qquad
 \sum_{j=1}^m d_j\leq C L^c.
 \label{eq:patch-cylinder-projection}
\end{equation}
All the summands in \eqref{eq:patch-cylinder-projection} have pairwise
orthogonal ranges as operators on \(\mathcal H_\Lambda\).
No bound on the internal tensor complexity of the individual
\(v_{j\ell}\) is asserted.
\end{proposition}

\begin{proof}
\emph{Nested squares and regularized filters.}
If \(u<4\), take \(P=\id\), expanded in an orthonormal basis on
\(X(u)\).  That set has at most \(81\) sites, so the expansion has
at most \(q^{81}\) terms.  This also covers empty square samples.
We henceforth suppose \(u\geq4\).

Choose
\[
 m=\lfloor u/2\rfloor+1,\qquad
 r_j=u+2(j-1),\qquad X_j=X(r_j),\qquad
 a_j=\frac{\kappa}{m},
\]
where \(0<\kappa\leq1/2\) will be fixed below using only the physical
parameters.  Then \(r_m\leq2u\), \(m\geq u/2\), and
\(\sum_j a_j=\kappa\).  The samples need not be distinct.

For \(R\geq0\) put \(b=e^{-R}\).  Minimize over density operators
\(x_j\) on \(\mathcal H_{X_j}\), independently for each \(j\), the
quantity
\begin{equation}
 \begin{split}
 N(R)&=\min_{\substack{x_j\geq0\\ \Tr x_j=1}}
               \norm{M(R,x)\Omega},\\
 M(R,x)&=L_m\cdots L_1,\qquad
 L_j=(x_j+b\id)^{-a_j/2}.
 \end{split}
 \label{eq:patch-variational-problem}
\end{equation}
Operators on \(X_j\) are extended by the identity on its complement.
The feasible set is compact, and every filter is positive and
invertible.  Since \(0\leq x_j\leq\id\),
\begin{equation}
 (1+b)^{-\kappa/2}
       \leq \norm{M(R,x)\Omega}
       \leq b^{-\kappa/2}.
 \label{eq:patch-elementary-norm-bounds}
\end{equation}
The lower bound follows by multiplying the least singular values of
the factors, and does not require them to commute.  Thus the minimum
is attained and is positive.  At any minimizing tuple set
\[
 \phi=\frac{L_m\cdots L_1\Omega}{N(R)},
 \qquad
 \rho_j=\Tr_{\Lambda\setminus X_j}
                \lvert\phi\rangle\langle\phi\rvert .
\]
All subsequent statements about minimizers apply to every minimizing
tuple, even when it is not unique.

The reason for this minimization is the following chain of estimates.
Stationarity will bound the excitation energy of $\phi$ independently
of $R$.  The gap will then make its marginals close enough to those of
$\Omega$ that the area law supplies a subspace of rank $e^{O(1+u)}$
carrying all but a fixed small fraction of each marginal.  This will force
the optimized norm to grow slowly:
$\log N(R)\le \kappa R/8+O(1+u)$.
To return from $\phi$ to $\Omega$, we will split each inverse filter
into a part of inside rank at most $e^R$ and a remaining part of norm at
most $(2e^{-R})^{a_j/2}$.  The term using only the remaining parts thus
has norm at most $N(R)(2e^{-R})^{\kappa/2}$, which decays exponentially
in $R$.  Every other term has a first finite-rank factor; nesting places
it in a cylinder at that factor's radius.  This is why the conclusion
uses several possible radii.

\emph{Descending stationarity without commuting filters.}
We first prove, in decreasing order of \(j\), that
\begin{equation}
 [L_j,\rho_j]=[x_j,\rho_j]=0.
 \label{eq:patch-stationarity-commutation}
\end{equation}
The trace identity behind the induction is the following.  If
\([L_k,\rho_k]=0\) and \(T\) is supported on \(X_k\), then
\begin{equation}
 \langle\phi,L_k T L_k^{-1}\phi\rangle
   =\Tr\bigl(L_k^{-1}\rho_k L_k T\bigr)
   =\langle\phi,T\phi\rangle .
 \label{eq:patch-trace-stripping}
\end{equation}
In a product of outer conjugations this identity can be applied
successively from the largest square inward: by nesting, the operator
remaining inside each conjugation is supported on that square.

Assume that \eqref{eq:patch-stationarity-commutation} has been proved
for \(k>j\), and let \(B\) be any anti-Hermitian operator on \(X_j\).
Vary \(x_j\) by \(x_j(s)=e^{sB}x_j e^{-sB}\), keeping all other
variables fixed.  Functional calculus gives
\[
 L_j(s)=e^{sB}L_j e^{-sB},\qquad
 \left.\frac{dL_j(s)}{ds}\right|_{s=0}L_j^{-1}
       =B-L_j B L_j^{-1}.
\]
For a differentiable invertible \(M(s)\), its logarithmic norm
derivative on the fixed input \(\Omega\) is
\[
 \left.\frac{d}{ds}\log\norm{M(s)\Omega}\right|_{s=0}
   =\operatorname{Re}
       \langle\phi,M'(0)M(0)^{-1}\phi\rangle .
\]
Here \(M'(0)M(0)^{-1}\) is \(B-L_j B L_j^{-1}\) conjugated
by \(L_m\cdots L_{j+1}\).  Strip those outer conjugations using
\eqref{eq:patch-trace-stripping}.  Stationarity at the minimum yields
\[
 0=\operatorname{Re}\Tr\bigl(\rho_j(B-L_j B L_j^{-1})\bigr)
   =-\operatorname{Re}\Tr\bigl(L_j^{-1}\rho_jL_j B\bigr).
\]
The term \(\Tr(\rho_jB)\) is purely imaginary.  An operator whose
real trace pairing with every anti-Hermitian operator vanishes is
Hermitian.  Consequently
\[
 L_j^{-1}\rho_jL_j=L_j\rho_jL_j^{-1},
 \qquad
 \rho_jL_j^2=L_j^2\rho_j .
\]
Positivity of \(L_j\) gives commutation with \(L_j\); functional
calculus then gives commutation with \(x_j\).  This starts at \(j=m\)
with no outer filters and proves the induction.  Notice that no
commutation between two different filters has been claimed.

In a common eigenbasis of \(x_j\) and \(\rho_j\), write \(x_i\)
and \(p_i\) for their eigenvalues.  A diagonal variation of \(x_j\)
within the probability simplex has logarithmic norm derivative
\begin{equation}
 -\frac{a_j}{2}\sum_i \frac{p_i}{x_i+b}\,dx_i.
 \label{eq:patch-simplex-derivative}
\end{equation}
Again, this follows by stripping the outer conjugations; the
diagonal variation commutes with \(x_j+b\id\).
For two indices with \(x_i,x_k>0\), mass can be moved in either
direction, so the ratios \(p_i/(x_i+b)\) agree.  Moving mass from
an index with \(x_i>0\) to any zero coordinate gives the one-sided
inequality.  Thus there is \(\lambda>0\) such that
\begin{equation}
 \frac{p_i}{x_i+b}=\lambda\quad(x_i>0),
 \qquad
 \frac{p_i}{b}\leq\lambda\quad(x_i=0).
 \label{eq:patch-kkt}
\end{equation}
Indeed, if the common ratio on the positive coordinates were zero,
the second inequality would force every \(p_i\) to be zero.
Writing \(r\) for the number of positive \(x_i\), we also have
\[
 \lambda(1+br)=\sum_{x_i>0}p_i\leq1,
 \qquad\hbox{so}\qquad 0<\lambda\leq1.
\]
Equivalently,
\begin{equation}
 x_i+b=\max\{p_i/\lambda,b\}.
 \label{eq:patch-clipped-eigenvalues}
\end{equation}
If \(\ell_i=(x_i+b)^{-a_j/2}\) is the corresponding eigenvalue
of \(L_j\), then for \(p_i,p_k>0\),
\begin{equation}
 \abs{\log(\ell_i/\ell_k)}
 \leq \frac{a_j}{2}\abs{\log(p_i/p_k)} .
 \label{eq:patch-log-ratio}
\end{equation}
This follows because \(t\mapsto\max\{t,\log b\}\) is
\(1\)-Lipschitz, applied to \(\log(p_i/\lambda)\).
The regularization keeps all \(\ell_i\) finite, including on zero
marginal eigenspaces; no ratio of zero probabilities is used.

\emph{A dimension-independent energy estimate.}
We next show that the minimizing vector has uniformly small
excitation energy.  The two distances
\(\abs{x-z}_\infty\) and \(\abs{y-z}_\infty\) at the endpoints
of a nearest-neighbor edge differ by at most one.  A closed square
of radius \(r\) cuts that edge only when
\[
 \min\{\abs{x-z}_\infty,\abs{y-z}_\infty\}
       \leq r
       <\max\{\abs{x-z}_\infty,\abs{y-z}_\infty\}.
\]
Since selected radii differ by at least two, an interaction crosses
at most one selected square.  If it crosses \(X_j\), its support
is disjoint from every \(X_k\) with \(k<j\), and is contained in
every \(X_k\) with \(k>j\).  This argument includes exact boundary
equalities and repeated clipped samples.

The ground equation gives
\[
 M H M^{-1}\phi=E_0\phi.
\]
For an interaction \(h\) crossing \(X_j\), remove the conjugations
by the containing outer filters using
\eqref{eq:patch-trace-stripping}.  The smaller filters commute
with \(h\) by disjointness.  Therefore
\begin{equation}
 \langle\phi,MhM^{-1}\phi\rangle
       =\langle\phi,L_j h L_j^{-1}\phi\rangle.
 \label{eq:patch-single-crossing}
\end{equation}
For a term crossing none of the selected squares the same argument
leaves its original expectation unchanged: strip the containing
filters, then commute past all disjoint ones.  On-site terms never
cross a square.

Consider now a Hermitian nearest-neighbor term \(h\) crossing
\(X_j\).  Use the Schmidt decomposition
\[
 \phi=\sum_{i:p_i>0}\sqrt{p_i}\,
          e_i\otimes f_i
\]
across \(X_j\) and its complement, where the \(e_i\) are chosen
in a simultaneous eigenbasis of \(L_j\) and \(\rho_j\).  Put
\[
 h_{ik}=\langle e_i\otimes f_i,\,
                       h(e_k\otimes f_k)\rangle.
\]
Hermiticity gives \(h_{ki}=\overline{h_{ik}}\).  Pairing the
\((i,k)\) and \((k,i)\) terms consequently gives
\begin{align}
 &\operatorname{Re}\langle\phi,L_j h L_j^{-1}\phi\rangle
           -\langle\phi,h\phi\rangle \notag\\
 &\qquad =
   \sum_{i,k:p_i p_k>0}
      \sqrt{p_i p_k}
      \bigl(\cosh(\log(\ell_i/\ell_k))-1\bigr)
      \operatorname{Re}h_{ik}.
 \label{eq:patch-hermitian-pairing}
\end{align}
For \(0\leq a\leq1\), the power series of the hyperbolic cosine
shows that
\[
 \cosh(ay)-1\leq a^2(\cosh y-1)\leq a^2\cosh y.
\]
Using \eqref{eq:patch-log-ratio} and
\(\sqrt{p_i p_k}\cosh(\tfrac12\log(p_i/p_k))
=(p_i+p_k)/2\), the nonnegative coefficient in
\eqref{eq:patch-hermitian-pairing} is bounded by
\begin{equation}
 \sqrt{p_i p_k}
      \bigl(\cosh(\log(\ell_i/\ell_k))-1\bigr)
       \leq \frac{a_j^2}{2}(p_i+p_k).
 \label{eq:patch-quadratic-coefficient}
\end{equation}
Zero Schmidt probabilities contribute zero to the expectation
and need no separate limiting argument.

Here is the step that prevents a Schmidt-rank factor from entering.
For operators \(A\) on \(X_j\) and \(B\) on its complement,
let \(A_{ik}=\langle e_i,Ae_k\rangle\) and
\(B_{ik}=\langle f_i,Bf_k\rangle\), restricted to the displayed
Schmidt vectors.  Cauchy--Schwarz in a row and in a column gives
\[
 \sum_k\abs{A_{ik}B_{ik}}\leq\norm{A}\norm{B},
 \qquad
 \sum_i\abs{A_{ik}B_{ik}}\leq\norm{A}\norm{B}.
\]
It follows that
\begin{equation}
 \sum_{i,k}(p_i+p_k)\abs{A_{ik}B_{ik}}
       \leq2\norm{A}\norm{B}.
 \label{eq:patch-weighted-row-column}
\end{equation}
Expand \(h\) in matrix units on its endpoint inside \(X_j\):
\[
 h=\sum_{a,b=1}^q
       \lvert a\rangle\langle b\rvert\otimes h_{ab},
 \qquad \norm{h_{ab}}\leq\norm{h}\leq J.
\]
The other factors are extended by identities within their respective
sides of the cut.  Combining
\eqref{eq:patch-hermitian-pairing}--\eqref{eq:patch-weighted-row-column}
term by term yields the explicit bound
\begin{equation}
 \abs{\operatorname{Re}\langle\phi,L_j h L_j^{-1}\phi\rangle
                  -\langle\phi,h\phi\rangle}
       \leq q^2J\,a_j^2.
 \label{eq:patch-one-term-energy}
\end{equation}

The ambient sample of a square of radius at most \(2u\) has at most
\(16u+4\) crossing lattice edges.  Intersecting with the physical
square cannot add a physical crossing edge.  Hence, using the
ground equation, \eqref{eq:patch-single-crossing}, and
\eqref{eq:patch-one-term-energy},
\begin{align}
 0\leq\langle\phi,(H-E_0\id)\phi\rangle
 &\leq q^2J(16u+4)\sum_{j=1}^m a_j^2 \notag\\
 &\leq 36q^2J\,\kappa^2.
 \label{eq:patch-energy-bound}
\end{align}
The last inequality uses \(m\geq u/2\) and \(u\geq4\).
This bound is independent of \(R\), the private dimensions of the
filters, and the choice of a minimizer.

Set \(C_E=36q^2J\), and now fix
\begin{equation}
 0<\kappa\leq
 \min\left\{\frac12,\sqrt{\frac{\Delta}{1024C_E}}\right\}.
 \label{eq:patch-kappa-choice}
\end{equation}
The full-system gap implies
\[
 1-\abs{\langle\Omega,\phi\rangle}^2
       \leq \frac{C_E\kappa^2}{\Delta}.
\]
In particular, up to a phase,
\(\norm{\phi-\Omega}\leq
\sqrt{2C_E/\Delta}\,\kappa\).
The trace norm distance of the two pure-state density operators is
\(2\sqrt{1-\abs{\langle\Omega,\phi\rangle}^2}\).
After partial trace, \eqref{eq:patch-kappa-choice} therefore gives
\begin{equation}
 \norm{\rho_j-\rho_{\Omega,X_j}}_1\leq\frac1{16}
 \quad\hbox{for every }j.
 \label{eq:patch-marginal-closeness}
\end{equation}
Only the original full-system gap has been used.

\emph{Constant-tail ranks from the area law.}
Proposition~\ref{prop:area-law}, together with the crossing-edge
count, gives a constant \(C_S\geq1\), depending only on the physical
parameters, such that
\[
 S_\Omega(X_j)\leq C_S(1+u).
\]
Let \(Q_j\) be the spectral projection of \(\rho_{\Omega,X_j}\)
onto eigenvalues at least
\(\exp[-16C_S(1+u)]\).  Summing the eigenvalues shows
\[
 \rank Q_j\leq e^{C_Q(1+u)},\qquad C_Q=16C_S.
\]
Markov's inequality for the marginal surprisal gives
\[
 \Tr\bigl(\rho_{\Omega,X_j}(\id-Q_j)\bigr)
       \leq\frac1{16}.
\]
By \eqref{eq:patch-marginal-closeness},
\begin{equation}
 \Tr\bigl(\rho_j(\id-Q_j)\bigr)\leq\frac18.
 \label{eq:patch-constant-tail}
\end{equation}
These \(Q_j\) come from the original ground state.  We have not
applied entropy continuity to the filtered state.

At a minimizer define the commuting positive product
\[
 T_j=\rho_j\,b(x_j+b\id)^{-1}.
\]
Equation~\eqref{eq:patch-kkt} shows that its eigenvalues obey
\[
 0\leq \frac{p_i b}{x_i+b}\leq b,
 \qquad
 \frac{p_i b}{x_i+b}\leq p_i.
\]
Thus \(0\leq T_j\leq b\id\) and \(T_j\leq\rho_j\).
The projector \(Q_j\) need not commute with \(T_j\); the positive
operator inequalities suffice to give
\begin{align}
 \Tr T_j
 &=\Tr(Q_jT_j)+\Tr((\id-Q_j)T_j) \notag\\
 &\leq b\,\rank Q_j+
          \Tr((\id-Q_j)\rho_j)
 \leq e^{-R+C_Q(1+u)}+\frac18.
 \label{eq:patch-regulator-trace}
\end{align}
Choose, for example,
\[
 C_0=C_Q+\log8,\qquad R_{\mathrm{start}}=C_0(1+u).
\]
Then \(\Tr T_j\leq1/4\) whenever \(R\geq R_{\mathrm{start}}\),
uniformly over all minimizers.

\emph{The nonsmooth minimum and its growth.}
The regulator estimate will bound the growth of $\log N(R)$.  Its rate
must stay below $\kappa/2$, the exponential decay rate of the
inverse-filter remainder described above.  Write
\[
 f(R,x)=\log\norm{M(R,x)\Omega},
 \qquad g(R)=\log N(R)=\min_x f(R,x).
\]
We justify the envelope calculation without choosing a continuous
family of minimizers.  For fixed \(x\),
\[
 \frac{\partial L_j}{\partial R}
       =\frac{a_j}{2}\,b(x_j+b\id)^{-1}L_j,
 \qquad
 \norm{\frac{\partial L_j}{\partial R}}
       \leq\frac{a_j}{2}\norm{L_j}.
\]
On \(0\leq R\leq T\), the product rule and
\eqref{eq:patch-elementary-norm-bounds} therefore give
\[
 \norm{\frac{\partial M}{\partial R}}
       \leq\frac{\kappa}{2}e^{\kappa T/2},
 \qquad
 \norm{M\Omega}\geq2^{-\kappa/2}.
\]
Consequently
\begin{equation}
 \abs{\partial_R f(R,x)}
       \leq \frac{\kappa}{2}\,2^{\kappa/2}e^{\kappa T/2}
       \quad(0\leq R\leq T).
 \label{eq:patch-uniform-local-lipschitz}
\end{equation}
The bound holds for every feasible \(x\) and is independent of all
Hilbert-space dimensions and of \(m\).  The minimum \(g\) inherits
the same Lipschitz bound on each bounded interval, and hence is
locally absolutely continuous.

Fix \(R\) and any minimizing tuple \(x_R\).  For \(h>0\),
\[
 g(R+h)-g(R)
       \leq f(R+h,x_R)-f(R,x_R).
\]
The upper right Dini derivative of \(g\) is therefore at most
\(\partial_R f(R,x_R)\).  At the minimizer the descending
trace rule removes the outer filters in each term of the product
derivative, giving
\[
 \partial_R f(R,x_R)
       =\sum_{j=1}^m\frac{a_j}{2}
          \Tr\bigl(\rho_j b(x_j+b\id)^{-1}\bigr).
\]
By \eqref{eq:patch-regulator-trace}, this is at most
\(\kappa/8\) for \(R\geq R_{\mathrm{start}}\).
At almost every point where \(g\) is differentiable its derivative
has this upper bound, so absolute continuity permits integration.
Using \(g(R_{\mathrm{start}})\leq
\kappa R_{\mathrm{start}}/2\) from
\eqref{eq:patch-elementary-norm-bounds}, we obtain the convenient
estimate
\begin{equation}
 g(R)\leq\frac{\kappa}{2}
                 \left(R_{\mathrm{start}}+\frac R4\right)
       \quad(R\geq R_{\mathrm{start}}).
 \label{eq:patch-minimum-growth}
\end{equation}
Neither differentiability of a minimizing tuple nor uniqueness of
that tuple is required.

\emph{An ordered expansion and a small all-tail remainder.}
Choose \(R\geq R_{\mathrm{start}}\) and a minimizing tuple at
that value.  Let \(\Pi_j\) be the spectral projection of \(x_j\)
onto eigenvalues at least \(b=e^{-R}\).  Split the inverse filter,
on its inside space, as
\[
 F_j=L_j^{-1},\qquad
 F_j^{\mathrm h}=F_j\Pi_j,\qquad
 F_j^{\mathrm t}=F_j(\id-\Pi_j).
\]
Since \(\Tr x_j=1\),
\begin{equation}
 \rank F_j^{\mathrm h}\leq b^{-1},
 \qquad
 \norm{F_j^{\mathrm t}}\leq(2b)^{a_j/2}.
 \label{eq:patch-head-tail-bounds}
\end{equation}
The rank in this display is on \(\mathcal H_{X_j}\), before
extension by the outside identity.

Inverting the ordered filter product gives
\(\Omega=N(R)F_1\cdots F_m\phi\).
Expand by the first head in this left-to-right order:
\begin{equation}
 F_1\cdots F_m
   =\sum_{j=1}^m
        F_1^{\mathrm t}\cdots F_{j-1}^{\mathrm t}
        F_j^{\mathrm h}F_{j+1}\cdots F_m
      +F_1^{\mathrm t}\cdots F_m^{\mathrm t}.
 \label{eq:patch-first-head-expansion}
\end{equation}
This identity uses distributivity only and never changes the order
of noncommuting factors.  Every factor preceding the \(j\)-th head
acts inside \(X_j\).  Hence the \(j\)-th summand applied to \(\phi\)
lies in
\[
 \mathcal C_j=S_j\otimes\mathcal H_{\Lambda\setminus X_j},
 \qquad
 \dim S_j\leq b^{-1},
\]
where \(S_j\) is the range on \(\mathcal H_{X_j}\) of
\(F_1^{\mathrm t}\cdots F_{j-1}^{\mathrm t}F_j^{\mathrm h}\).
The preceding factors may alter the head subspace, but cannot
increase its dimension.

The remaining vector has norm at most
\[
 N(R)\prod_j\norm{F_j^{\mathrm t}}
       \leq N(R)(2b)^{\kappa/2}.
\]
By \eqref{eq:patch-minimum-growth}, its logarithm is at most
\begin{equation}
 \frac{\kappa}{2}
       \left[C_0(1+u)+\log2-\frac{3R}{4}\right].
 \label{eq:patch-all-tail-log}
\end{equation}
To make this choice explicit and uniform, set
\begin{equation}
 R_0=\frac43\left[
           C_0\left(B_0+\frac1{\log2}\right)
                 +1+\frac{2p}{\kappa}\right],
 \qquad R=R_0\log L.
 \label{eq:patch-explicit-radius}
\end{equation}
For \(L\geq2\) and \(u\leq B_0\log L\),
\[
 C_0(1+u)\leq
       C_0\left(B_0+\frac1{\log2}\right)\log L,
 \qquad \log2\leq\log L.
\]
Thus \(R\geq R_{\mathrm{start}}\), and
\eqref{eq:patch-all-tail-log} is at most \(-p\log L\).
Writing \(\mathcal E=\sum_j\mathcal C_j\), we have consequently
proved
\begin{equation}
 \dist(\Omega,\mathcal E)\leq L^{-p},
 \qquad
 \dim S_j\leq L^{R_0}.
 \label{eq:patch-distance-to-cylinders}
\end{equation}

\emph{Orthogonalization at variable radii.}
Let \(P\) be the orthogonal projection onto \(\mathcal E\).
We finish by putting it in the required form, keeping the radius
associated with each new summand.  Process the cylinders in
increasing \(j\).  By nesting, the earlier span can be written
\[
 \mathcal E_{j-1}:=\sum_{i<j}\mathcal C_i
       =W_j\otimes\mathcal H_{\Lambda\setminus X_j}
\]
for a subspace \(W_j\subseteq\mathcal H_{X_j}\).
Let \(P_{W_j}\) be its orthogonal projection, and define
\[
 \widetilde S_j=(\id-P_{W_j})S_j.
\]
This is the \emph{projected image} of \(S_j\), not its intersection
with \(W_j^\perp\).  It has dimension at most \(\dim S_j\), and
\[
 \mathcal E_{j-1}+\mathcal C_j
    =\mathcal E_{j-1}\ \mathbin{\oplus}\
       \bigl(\widetilde S_j
                    \otimes\mathcal H_{\Lambda\setminus X_j}\bigr).
\]
Choose an orthonormal basis \(v_{j\ell}\) of each
\(\widetilde S_j\); if that space is zero there are no terms for
that \(j\).  The resulting cylinder ranges are pairwise orthogonal
and their direct sum is \(\mathcal E\), proving
\eqref{eq:patch-cylinder-projection}.  Because \(X_j\subseteq X(2u)\),
this projection is supported on \(X(2u)\).  Equation
\eqref{eq:patch-distance-to-cylinders} gives its stated error.

Finally,
\[
 \sum_j d_j\leq mL^{R_0}
       \leq \left(\frac{B_0}{2}\log L+1\right)L^{R_0}.
\]
This is bounded by \(C L^{R_0+1}\); enlarging \(C\) covers the
previous bounded-\(u\) case.  The choices of \(\kappa,C_0,R_0\)
depend only on the fixed parameters in the proposition, and all
bounds are uniform in the center, the Hamiltonian, and \(L\).
\end{proof}
